\documentclass[12pt]{article}

\usepackage{graphicx}
\usepackage{amsmath}
\usepackage{booktabs}
\usepackage[a4paper, margin=1in]{geometry}
\usepackage{hyperref}

\title{Diabetes Risk Prediction and Guideline-Grounded Recommendation System}
\author{
\begin{tabular}{ccc}
Shubham Aggarwal & Atharva Sharma & Bhavya Punj \\
\end{tabular}
\\[0.5cm]
\textit{Final Project}
}
\date{\today}

\begin{document}

\maketitle

\begin{abstract}
This project combines a deployed logistic-regression screening model with a guideline-grounded retrieval and language-generation layer. The Streamlit application collects patient health indicators, computes a diabetes risk probability from the shipped model artifacts, extracts salient risk factors, retrieves relevant passages from medical guideline PDFs using FAISS and a CrossEncoder reranker, and produces structured JSON guidance through Gemini with a safe fallback when the API key is unavailable.
\end{abstract}

\section{Project Overview}

Early detection of diabetes is important because screening systems can help flag elevated risk before complications develop. The current project is not only a classifier; it is a complete screening and guidance workflow that takes user inputs, predicts risk, explains the main factors, recommends appropriate specialists, and grounds the final recommendations in guideline passages.

The runtime application is the Streamlit interface in \texttt{app/streamlit\_app.py}, with a thin launcher in \texttt{Final/streamlit\_app.py}. The app loads a pretrained logistic-regression model and scaler from \texttt{models/lr\_model.pkl} and \texttt{models/scaler.pkl}, then uses a LangGraph pipeline for the downstream recommendation flow.

\section{Dataset and Feature Set}

The model was trained from the BRFSS 2015 diabetes health indicators dataset, available on Kaggle as the Diabetes Health Indicators Dataset.

The exported training script in \texttt{Research/Final-(model)/export\_model.py} reads \texttt{Research/data/dataset\_binary.csv}, removes redundant columns, standardizes the selected inputs, optionally balances the classes with SMOTE, and fits a logistic-regression classifier.

\textbf{Data characteristics:}
\begin{itemize}
\item Approximately 253,680 records
\item Binary target: \texttt{Diabetes\_binary}
\item Feature set used by the deployed classifier: 19 variables
\item Removed during export: \texttt{Stroke} and \texttt{HeartDiseaseorAttack}
\end{itemize}

The exact model feature order is stored in \texttt{models/model\_metadata.json}:

\begin{itemize}
\item \texttt{HighBP}, \texttt{HighChol}, \texttt{CholCheck}, \texttt{BMI}
\item \texttt{Smoker}, \texttt{PhysActivity}, \texttt{Fruits}, \texttt{Veggies}
\item \texttt{HvyAlcoholConsump}, \texttt{AnyHealthcare}, \texttt{NoDocbcCost}
\item \texttt{GenHlth}, \texttt{MentHlth}, \texttt{PhysHlth}, \texttt{DiffWalk}
\item \texttt{Sex}, \texttt{Age}, \texttt{Education}, \texttt{Income}
\end{itemize}

The UI also collects glucose as a user-facing input for risk-factor analysis and guideline retrieval, but it is not part of the classifier feature vector. This separation is intentional and reflected in \texttt{app/streamlit\_app.py}.

\section{Model Development and Deployment}

The training/export path is intentionally simple and reproducible. The final artifact set contains:

\begin{itemize}
\item \texttt{lr\_model.pkl}
\item \texttt{scaler.pkl}
\item \texttt{model\_metadata.json}
\end{itemize}

The export script standardizes the selected features with \texttt{StandardScaler}, applies SMOTE when available, and fits \texttt{LogisticRegression(C=0.01, max\_iter=1000, random\_state=42)}. The deployed decision threshold is 0.58, which is also recorded in the metadata file and used by the Streamlit app.

The app loads the model artifacts at runtime and computes the probability of the positive class. The displayed risk label is derived from the threshold rather than from a hardcoded class prediction, which makes the UI easier to interpret during screening.

\subsection*{Why Logistic Regression}

\begin{itemize}
\item It is transparent and easy to audit for a healthcare screening use case.
\item It supports stable probability outputs for threshold-based decision making.
\item It is lightweight enough for a Streamlit deployment with cached artifacts.
\end{itemize}

\subsection*{Feature Handling}

\begin{itemize}
\item Categorical fields such as education and income are mapped to ordinal codes.
\item Binary health indicators are converted to 0/1 features.
\item Age is converted into the category scheme used by the model.
\item The model input order is fixed to match the exported artifact metadata.
\end{itemize}

\section{Research Benchmark and Model Selection Rationale}

Before freezing the deployment pipeline, we benchmarked multiple classifiers during the research phase and compared their screening-oriented performance. These benchmark results motivated the final choice of logistic regression for production use.

\subsection*{Benchmark Models}

\begin{itemize}
\item Logistic Regression (LR)
\item Random Forest (RF)
\item XGBoost (XGB)
\item Artificial Neural Network (ANN)
\end{itemize}

\subsection*{Benchmark Metrics (Research Phase)}

\begin{table}[h]
\centering
\begin{tabular}{lcccc}
\toprule
Model & Accuracy & Precision & Recall & F1 \\
\midrule
Logistic Regression & 0.765 & 0.356 & 0.654 & 0.461 \\
Random Forest & 0.779 & 0.369 & 0.613 & 0.460 \\
XGBoost & 0.786 & 0.378 & 0.609 & 0.466 \\
ANN & 0.782 & 0.374 & 0.622 & 0.467 \\
\bottomrule
\end{tabular}
\caption{Research-phase benchmark comparison across candidate models}
\end{table}

These values are retained from the milestone benchmarking stage and are included here to document the model selection process.

\subsection*{Confusion Matrix Comparison}

\begin{figure}[h]
\centering
\includegraphics[width=0.48\textwidth]{cm_lr.png}
\caption{Confusion Matrix for Logistic Regression (research phase)}
\end{figure}

\begin{figure}[h]
\centering
\includegraphics[width=0.48\textwidth]{cm_ann.png}
\caption{Confusion Matrix for ANN (research phase)}
\end{figure}

\subsection*{Why Logistic Regression Was Deployed}

Although ANN and XGBoost showed slightly better aggregate metrics in some benchmark views, logistic regression was selected for deployment because it offered a favorable trade-off for healthcare screening:

\begin{itemize}
\item stronger recall profile for identifying at-risk users during the benchmark phase,
\item lower false-negative tendency in screening-oriented interpretation,
\item simpler and more transparent behavior for auditability,
\item lightweight inference for reliable Streamlit deployment.
\end{itemize}

In screening scenarios, missing truly at-risk patients is often costlier than raising additional follow-up checks. This practical consideration is central to the final deployment decision.

\section{User Interface and Prediction Flow}

The Streamlit interface is the user entry point. It is split into lifestyle, medical history, health condition, and demographic sections to make data entry understandable for non-technical users. After form submission, the app stores the raw inputs, transforms the classifier features, and then routes the request into the AI workflow.

The end-to-end flow is:

\begin{itemize}
\item collect user inputs in the Streamlit form
\item transform the feature vector and compute risk probability
\item extract notable risk factors from the raw inputs
\item recommend suitable specialists
\item retrieve and rerank guideline passages
\item build a structured prompt
\item generate a JSON response through Gemini or a fallback response
\end{itemize}

The app also includes a short debug trace mode, cached artifact loading, and a graceful fallback if the Gemini API key is missing or the model artifacts are unavailable.

\section{Risk Factor and Specialist Logic}

The risk-factor extraction logic is implemented in \texttt{agent/utils.py}. It looks for clinically relevant patterns in the user inputs, including:

\begin{itemize}
\item elevated or borderline glucose
\item high BMI or overweight status
\item high blood pressure and high cholesterol
\item smoking, low activity, low fruit and vegetable intake, and heavy alcohol use
\item age-related risk
\end{itemize}

If no major signal is found, the system returns a neutral explanation rather than fabricating a risk factor.

Specialist recommendations are generated in \texttt{agent/doctor.py}. The logic is rule-based and favors endocrinology for higher probabilities, cardiology for blood-pressure-related concern, nutrition support for high BMI, and a general physician as the default fallback.

\section{Retrieval-Augmented Generation}

The guidance layer is built around two guideline PDFs in \texttt{data/guidelines}:

\begin{itemize}
\item \texttt{who\_diabetes\_diagnosis.pdf}
\item \texttt{idf\_diabetes\_management.pdf}
\end{itemize}

The second local file (\texttt{idf\_diabetes\_management.pdf}) is sourced from NICE guideline NG28: Type 2 diabetes in adults: management.

The retrieval pipeline is implemented in \texttt{agent/rag\_faiss.py} and works as follows:

\begin{itemize}
\item PDF text is extracted with \texttt{PyPDF2}
\item Text is chunked into 150-word windows with 30-word overlap
\item Chunks are embedded using \texttt{all-MiniLM-L6-v2}
\item A FAISS L2 index stores and searches the embeddings
\item Candidate passages are reranked with \texttt{cross-encoder/ms-marco-MiniLM-L-6-v2}
\end{itemize}

The app then combines the reranked passages into a compact context string for prompt construction. This keeps the generated response anchored to retrieved medical text rather than relying on the language model alone.

\subsection*{LangGraph Orchestration}

The workflow graph in \texttt{agent/langgraph\_flow.py} is linear and explicit:

\begin{center}
ML Prediction $\rightarrow$ Risk Factors $\rightarrow$ Specialist Recommendation $\rightarrow$ Retrieval $\rightarrow$ LLM Response
\end{center}

This structure makes each stage inspectable and easy to debug. The higher-level wrapper in \texttt{agent/workflow.py} also normalizes source names for presentation in the UI.

\section{LLM Response Format}

The prompt builder in \texttt{agent/prompt.py} instructs Gemini to return a strict JSON object containing:

\begin{itemize}
\item an explanation of the risk result
\item practical recommendations
\item preventive measures
\item suggested specialists
\item source citations
\item a screening-only disclaimer
\end{itemize}

The Gemini wrapper in \texttt{agent/llm.py} uses \texttt{gemini-2.5-flash} when \texttt{GEMINI\_API\_KEY} is available. If the API key is missing or the request fails, the system returns a deterministic fallback JSON payload so the app can still render a complete result page.

The response parser in the Streamlit app validates the JSON and degrades safely if the returned text is malformed.

\section{Implementation Notes}

\begin{itemize}
\item The deployed application uses cached artifact loading to avoid repeated model deserialization.
\item The UI presents a risk percentage and a high/low label computed from the 0.58 threshold.
\item The source labels shown to users are normalized for readability in the result panel.
\item The system is designed as a screening aid and explicitly avoids presenting itself as a diagnostic tool.
\end{itemize}

\section{Limitations and Future Work}

The current system is intentionally constrained to a small, well-defined guideline corpus and a single deployed classifier. That makes the pipeline easier to audit, but it also limits generalization to broader medical scenarios.

Potential next steps include:

\begin{itemize}
\item adding explainability outputs such as SHAP-based feature attributions
\item expanding the RAG corpus with additional trusted clinical sources
\item introducing evaluation logging for the deployed screening workflow
\item formalizing calibration and threshold analysis for different risk tolerances
\end{itemize}

\section{Conclusion}

This project delivers a complete diabetes risk screening and guidance application rather than a standalone classifier. The deployed logistic-regression model provides the risk score, the retrieval layer grounds the advice in guideline text, and the Gemini-backed JSON response packages the result into a user-facing recommendation panel. The architecture is lightweight, reproducible, and aligned with a healthcare screening workflow.

\section{Project Resources}

\begin{itemize}
\item Live app: \href{https://diabetespreidictor.streamlit.app/}{Diabetes Risk Predictor}
\item Runtime entry point: \texttt{app/streamlit\_app.py}
\item Model export script: \texttt{Research/Final-(model)/export\_model.py}
\item RAG implementation: \texttt{agent/rag\_faiss.py}
\end{itemize}

\begin{thebibliography}{9}

\bibitem{kaggle_dataset}
A. Teboul,
\textit{Diabetes Health Indicators Dataset}, Kaggle.
\href{https://www.kaggle.com/datasets/alexteboul/diabetes-health-indicators-dataset}{Link}

\bibitem{who}
World Health Organization,
\textit{Classification of Diabetes Mellitus}, 2019.
\href{https://apps.who.int/iris/bitstream/handle/10665/325182/9789241515702-eng.pdf}{Link}

\bibitem{guidelines}
National Institute for Health and Care Excellence (NICE),
\textit{Type 2 diabetes in adults: management (NG28)}.
\href{https://www.nice.org.uk/guidance/ng28/resources/type-2-diabetes-in-adults-management-pdf-1837338615493}{Download guidance (PDF)}

\end{thebibliography}

\end{document}